% パッケージの読み込み
\usepackage{graphicx}
\usepackage{mathtools} % 含まれる amsmath
\usepackage{amssymb}
\usepackage{amsfonts}
\usepackage{amsthm}
\usepackage{bm}
\usepackage{mleftright}
\usepackage{optidef}

\usepackage{tikz}
\usetikzlibrary{shapes,positioning}
\usepackage{pgfplots}
\pgfplotsset{compat=1.18}
\usepackage{framed,color}

\usepackage{listings}

\usepackage{fancyhdr}
\usepackage{multicol}
\usepackage{comment}
\usepackage{booktabs}
\usepackage{ascmac}
\usepackage{url}
\usepackage{lscape}
\usepackage[normalem]{ulem}
\usepackage[hypertexnames=false]{hyperref}
\hypersetup{% hyperrefオプションリスト
 bookmarksopen=true,%
 colorlinks=true,%
 linkcolor=blue,
 citecolor=blue,
}
\usepackage{cleveref}
%\usepackage{autonum}
\usepackage[numbers,sort]{natbib}

% アルゴリズム関連のパッケージ
\usepackage{algorithmicx}
\usepackage{algpseudocode}

% カラーの定義
\definecolor{shadecolor}{gray}{0.90}

% cleveref の設定
\crefformat{chapter}{第#2#1#3章}
\crefformat{section}{#2#1#3節}
\crefformat{subsection}{#2#1#3節}
\crefname{figure}{図}{図}
\crefname{table}{表}{表}
\crefname{theorem}{定理}{定理}
\crefname{equation}{式}{式}
\crefname{appendix}{付録}{付録}
\newcommand{\crefrangeconjunction}{--}
\newcommand{\crefpairconjunction}{, }
\newcommand{\crefmiddleconjunction}{, }
\newcommand{\creflastconjunction}{, }

% 定理環境の設定
\theoremstyle{definition}
\renewcommand{\proofname}{\normalfont{\uline{証明}}}

% ペアとなる区切り記号の定義
\DeclarePairedDelimiter{\paren}{\lparen}{\rparen}
\DeclarePairedDelimiter{\tuple}{\lparen}{\rparen}
\DeclarePairedDelimiter{\sqbra}{\lbrack}{\rbrack}
\DeclarePairedDelimiter{\set}{\lbrace}{\rbrace}
\DeclarePairedDelimiter{\abs}{\lvert}{\rvert}
\DeclarePairedDelimiter{\card}{\lvert}{\rvert}
\DeclarePairedDelimiter{\norm}{\lVert}{\rVert}
\DeclarePairedDelimiterX{\pgiven}[2]{\lparen}{\rparen}{#1\,\delimsize\vert\,#2}
\DeclarePairedDelimiterX{\qgiven}[2]{\lbrack}{\rbrack}{#1\,\delimsize\vert\,#2}

% 数学演算子の定義
\DeclareMathOperator{\rank}{rank}
\DeclareMathOperator{\Int}{int}
\DeclareMathOperator{\ke}{Ker}
\DeclareMathOperator*{\minimize}{minimize}
\DeclareMathOperator*{\argmin}{arg\,min}
\DeclareMathOperator*{\argmax}{arg\,max}
\DeclareMathOperator{\Tr}{Tr}
\DeclareMathOperator{\Diag}{Diag}
\DeclareMathOperator{\subjectto}{subject~to}

% アルゴリズム関連の設定
\newcommand{\Step}[1]{\algrenewcommand{\alglinenumber}[1]{\textbf{Step ##1:} } #1}



% 左線付き環境の定義 (needs framed, color)
\newenvironment{myleftbar}[1]{
  \def\FrameCommand{{\color{#1}\vrule width 1.5pt }\hspace{10pt}}
  \MakeFramed {\advance\hsize-\width \FrameRestore}}
{\endMakeFramed\par}

% 定理環境の基礎となるカウンターを定義
\newtheorem{mainthm}{定理}


% 定理環境の定義（左線付き）

% その他の定理環境の定義（共通の番号付け）
\theoremstyle{definition}
\newtheorem{assumption}[mainthm]{仮定}
\newtheorem{theorem}{定理}
\newtheorem{proposition}[mainthm]{命題}
\newtheorem{definition}[mainthm]{定義}
\newtheorem{lemma}[mainthm]{補題}
\newtheorem{remark}[mainthm]{注意}
\newtheorem{corollary}[mainthm]{系}
\newtheorem{example}[mainthm]{例}
\newtheorem{conjecture}[mainthm]{予想}

% tcolorbox の環境設定





% ベクトルの表記
\newcommand{\bvec}[1]{\boldsymbol{#1}}

% その他のコマンド
\newcommand{\relmiddle}[1]{\mathrel{}\middle#1\mathrel{}}
\newcommand{\st}{\mathrm{s.t.}}
